The account calls the Rust \code{lower} and \code{lower\_scoped} on install and
\code{rule\_eval::}\allowbreak\code{authorize} in \code{\_\_check\_auth}.
Theorem~\ref{stmt:lowering-equivalence} does for lowering what
Theorem~\ref{stmt:total-equivalence-of-the-evaluator} does for \code{authorize};
Theorem~\ref{stmt:canonical-identity}, in Perch's own Lean package, is about a
grant's \code{doc\_hash}, the SHA-256 of its CANON~v1 bytes (RFC~8785~\cite{rfc8785}
restricted).

% Generated from source/statements.json.
\begin{perchstatement}{12}{Lowering equivalence.}\label{stmt:lowering-equivalence}
\[\begin{aligned}&\mathit{LowerRoom}(d) \implies \exists R.\ \mathit{lower}_{\mathrm{T}}(d, A) = \mathsf{ok}(R)\\
&\qquad \land\ \llbracket R \rrbracket = \mathit{lower}(\llbracket d \rrbracket, \llbracket A \rrbracket)\\
&\mathit{ArgsRoom}(d) \implies \exists y.\\
&\qquad \mathit{lowerScoped}_{\mathrm{T}}(d, A, X, g) = \mathsf{ok}(y)\\
&\qquad \land\ \llbracket y \rrbracket = \mathit{lowerScoped}(\llbracket d \rrbracket,\\
&\qquad\qquad \llbracket A \rrbracket, \llbracket X \rrbracket, \llbracket g \rrbracket)\end{aligned}\]
Translated \code{lower} returns the model's lowering when the document leaves room in its bounded vectors (\code{LowerRoom}), and translated \code{lower\_scoped} does, refusals included, under the weaker \code{ArgsRoom}. For a valid document the lowered rules allow exactly what the document allows plus, when recovery is configured, every \code{apply\_doc} call through the recovery rule, which the shell refuses outside a completion (tested; O8).
\statementhypotheses{\code{LowerRoom} for \code{lower}, \code{ArgsRoom} for \code{lower\_scoped}; \code{Doc.Valid} for the semantic part}
\statementevidence{\repolink{https://github.com/bitswell/soroban-rule-core/blob/cbabe594ae2df408991e825e56ddd618ea29541b/formal/RuleCore/Equiv/Lowering.lean\#L602}{Equiv.lower\_eq}; \repolink{https://github.com/bitswell/soroban-rule-core/blob/cbabe594ae2df408991e825e56ddd618ea29541b/formal/RuleCore/Equiv/Lowering.lean\#L717}{lower\_scoped\_eq} \ensuremath{\cdot} \statusobj{PROVED \ensuremath{\cdot} TRANSLATION}; \repolink{https://github.com/bitswell/soroban-rule-core/blob/cbabe594ae2df408991e825e56ddd618ea29541b/formal/RuleCore/Theorems/ExtractedLowering.lean\#L53}{Extracted.lower\_allows}; \repolink{https://github.com/bitswell/soroban-rule-core/blob/cbabe594ae2df408991e825e56ddd618ea29541b/formal/RuleCore/Theorems/ExtractedLowering.lean\#L84}{lower\_authorize} \ensuremath{\cdot} \statusobj{PROVED \ensuremath{\cdot} MODEL + TRANSLATION}}
\end{perchstatement}

% Generated from source/statements.json.
\begin{perchstatement}{13}{Canonical identity.}\label{stmt:canonical-identity}
\[\mathit{emitDoc}(d_1) = \mathit{emitDoc}(d_2) \implies d_1 = d_2\]
Perch's canonical encoding is injective on the modeled document type; hash identity then holds up to SHA-256 collisions, an assumption outside the theorem. The theorem is about a Lean emitter written to match Perch's Rust \code{canon.rs} and linked to it by golden fixtures and fuzzing, not extraction; JSON parsing and compiler correctness need separate assurance, and no theorem connects this encoding to rule-core's document type. The CANON version is not part of the hashed bytes, so a \code{doc\_hash} is read under the identity the account reports (\code{capabilities().\allowbreak{}doc\_\allowbreak{}identity}); three domain separations are proven: \code{configPreimage\_\allowbreak{}ne\_\allowbreak{}emitDoc}, \code{rulePreimage\_\allowbreak{}ne\_\allowbreak{}configPreimage} and \code{rulePreimage\_\allowbreak{}ne\_\allowbreak{}emitDoc} (rule, recovery-configuration and document preimages are pairwise distinct).
\statementhypotheses{The modeled document type}
\statementevidence{\repolink{https://github.com/bitswell-hq/stellar-registry__perch/blob/0b3752d8a179a5952b05fbebd29131dbbf99fa80/formal/PerchFormal/CanonProofs.lean\#L681}{PerchFormal.Canon.emitDoc\_injective}; \repolink{https://github.com/bitswell-hq/stellar-registry__perch/blob/0b3752d8a179a5952b05fbebd29131dbbf99fa80/formal/PerchFormal/CanonProofs.lean\#L694}{canonical\_bytes\_identity} \ensuremath{\cdot} \statusobj{PROVED \ensuremath{\cdot} PERCH CANON PACKAGE}; domain separation: \repolink{https://github.com/bitswell-hq/stellar-registry__perch/blob/0b3752d8a179a5952b05fbebd29131dbbf99fa80/formal/PerchFormal/CanonProofs.lean\#L722}{configPreimage\_ne\_emitDoc}, \repolink{https://github.com/bitswell-hq/stellar-registry__perch/blob/0b3752d8a179a5952b05fbebd29131dbbf99fa80/formal/PerchFormal/CanonProofs.lean\#L760}{rulePreimage\_ne\_configPreimage}, \repolink{https://github.com/bitswell-hq/stellar-registry__perch/blob/0b3752d8a179a5952b05fbebd29131dbbf99fa80/formal/PerchFormal/CanonProofs.lean\#L765}{rulePreimage\_ne\_emitDoc}}
\end{perchstatement}


\noindent Perch's
integration suite runs rule-core's \fact{lean-doc-traces} Lean-generated document
traces through natively compiled \code{apply\_doc}, \code{apply\_scope} and
\code{check\_auth}, with the real Ed25519 verifier, enforcing authorization and
the generation set from each trace, and checks the stored rules, every verdict
and every spending record against the model; Perch's CI also replays rule-core's
conformance vectors. The theorems hold over one fixed rule slice, while the shell
builds a slice of the selected rules for each authorization; carrying them
through that selection is obligation O2 (Table~\ref{tab:tab-obligations}).

\perchtable{tab-obligations}

% Generated from source/statements.json.
\begin{perchproposition}{1}{From the evaluator to the account.}\label{stmt:from-the-evaluator-to-the-account}
Under obligations O1--O8, Theorem~\ref{stmt:total-equivalence-of-the-evaluator} with \code{RuleSet.\allowbreak{}authorize\_\allowbreak{}chosen}, the per-rule forms of Theorem~\ref{stmt:grant-target-containment} and Theorem~\ref{stmt:the-containment-chain} (\code{Doc.\allowbreak{}lowerScoped\_\allowbreak{}contract}, \code{Doc.\allowbreak{}lowerScoped\_\allowbreak{}allows}), and \code{authorize\_\allowbreak{}generation} (Theorem~\ref{stmt:recovery-invalidation}) give, for each context whose selected rule belongs to a grant: the context is allowed by the grant's source document, calls neither the account nor its excluded set, and is refused after a completed recovery.
\statementhypotheses{Obligations O1--O8. O1, O3, O4 and O6--O8 are tested, O2 is argued and tested, and O5, that compiler outputs satisfy \code{Doc.Valid}, is unproven (\code{\fact{claim-perch-validation}}); every grant rule satisfies \code{ArgsRoom} (assumed, \code{rule.lower-room}); the account's Wasm, built by rustc \fact{rustc-perch}, computes what the translation computes (trusted; nothing checks it)}
\statementevidence{Argued from the theorems above and the obligations; not mechanized \ensuremath{\cdot} \statusobj{ARGUED}}
\end{perchproposition}


\paragraph{What a signature binds.} A signature must commit to the rule IDs the
evaluator decides over. With $H$ = SHA-256, the account's
\fact{contract-id-bytes}-byte ID $\mathit{id}_A$, the tag $\tau$ =
\facthash{auth-domain-tag} and the selected rule IDs $\mathit{ids}$, an external
signer signs $D$:
\begin{align*}
  P &= H\bigl(\mathrm{XDR}(\mathtt{SorobanAuthorization}\{\mathit{nid}, \mathit{nonce}, \mathit{exp}, \mathit{inv}\})\bigr),\\
  B &= H(\tau \,\|\, \mathit{id}_A \,\|\, P), \qquad
  D = H\bigl(B \,\|\, \mathrm{XDR}(\mathtt{Vec}\langle\mathtt{u32}\rangle\ \mathit{ids})\bigr).
\end{align*}
$P$ is the host's payload for a legacy \code{Address} credential; CAP-0071's
address-bound credentials~\cite{cap-0071} hash \code{SorobanAuthorizationWithAddress}, which adds $A$. A
delegated signer approves the host's address-bound payload, which omits
$\mathit{ids}$. In OpenZeppelin's account~\cite{oz-stellar-accounts} at \fact{oz-v080-tag}, an
external signer signs $H(\mathrm{XDR}(\mathtt{AuthDigestPreimage}\{A, P,
\mathit{ids}\}))$ and a delegated one authorizes that preimage through \code{require\_auth\_for\_args}
(Table~\ref{tab:tab-binding}). The lemma is a paper proof, not mechanized.

\begin{lemma}[binding]\label{lem:binding}
Suppose the verifier named for an external signer correctly verifies an
EUF-CMA-secure scheme and $H$ is collision resistant. An adversary given
signatures on $D(t_1),\dots,D(t_q)$ for tuples
$t_i = (A, \mathit{nid}, \mathit{nonce}, \mathit{exp}, \mathit{inv}, \mathit{ids})$
of its choice outputs a $t \notin \{t_i\}$ with a signature accepted on $D(t)$
only with negligible probability.
\end{lemma}

\begin{proof}
The map from $t$ to $D$ composes injective encodings with $H$: XDR encodes
$(\mathit{nid}, \mathit{nonce}, \mathit{exp}, \mathit{inv})$ injectively, and its
union discriminant separates the two credential forms; $\tau$ is constant and
$\mathit{id}_A$ and $P$ have fixed lengths; $B$ has a fixed length and the XDR of $\mathit{ids}$ is
self-delimiting. If $D(t) = D(t_i)$ for some $i$, this gives a collision of $H$;
otherwise the output is a forgery. For WebAuthn the signed message carries $D$ as
the client-data challenge. \qed
\end{proof}

\noindent The lemma covers signatures on $D$ values only; a key that also signs
elsewhere, such as a Stellar account key, is covered only if no message it signs
there can equal a $D$ value. $D$ does not depend on the contents of the rules
$\mathit{ids}$ names, the account's revision, the transaction's other entries or
its fee payer.

\paragraph{Domain tags.} Perch separates its hash domains with
\facthash{auth-domain-tag} (the payload $B$, scheme
\fact{auth-scheme-name}), \fact{domain-tag-rule} (one rule's bytes),
\fact{domain-tags-recovery} (recovery statements, configurations, credentials and
replacement sets) and, as Poseidon2 constants, \fact{domain-tags-zk}. CANON
document hashes carry no tag; \code{configPreimage\_ne\_emitDoc},
\code{rulePreimage\_ne\_configPreimage} and \code{rulePreimage\_ne\_emitDoc} in
Perch's canonical-encoding package prove that configuration, rule and document
preimages never coincide. Each SHA-256 tag begins its preimage and none is a
prefix of another, so preimages from different domains differ (by inspection,
not proved). oz-policy-builder uses \fact{opb-domain-tags}.
